\documentclass[11pt,a4paper]{amsart}

\usepackage[margin=1in]{geometry}
\usepackage{amsmath,amsthm,amssymb,mathtools}

\DeclareMathOperator{\rad}{rad}
\newcommand{\N}{\mathbb N}

\newtheorem{theorem}{Theorem}[section]
\newtheorem{lemma}[theorem]{Lemma}
\newtheorem{proposition}[theorem]{Proposition}

\title{Totient Fibre Extremes}
\author{GPT-5.5 Pro}

\begin{document}

\begin{abstract}
Let \(f_{\max}(n)\) and \(f_{\min}(n)\) denote respectively the largest and smallest integers \(m\) with \(\phi(m)=n\), whenever such integers exist. We determine the maximal possible separation between the largest and smallest elements of a totient fibre. More precisely, if the maximum is taken over totient values \(n\le x\), then
\[
\max_{\substack{n\le x\\ n\in \phi(\N)}}\frac{f_{\max}(n)}{f_{\min}(n)}
=(e^\gamma+o(1))\log\log x .
\]
The upper bound follows from the extremal order of \(m/\phi(m)\), while the lower bound is obtained by a construction using Linnik's theorem on the least prime in an arithmetic progression. We also record a simple permanence observation showing that any one nontrivial totient collision propagates to infinitely many.
\end{abstract}

\maketitle

\section{Preliminaries}

Throughout, \(\phi\) denotes Euler's totient function and \(\gamma\) denotes Euler's constant. Since \(f_{\min}(n)\) and \(f_{\max}(n)\) are defined only when \(n\) is a totient value, we interpret the expression in the problem as
\[
\mathcal R(x):=
\max_{\substack{n\le x\\ n\in \phi(\N)}}
\frac{f_{\max}(n)}{f_{\min}(n)} .
\]
The maximum is nonempty for \(x\ge 1\), since \(1=\phi(1)=\phi(2)\).

We first note that every totient fibre is finite. Indeed, for an odd prime power \(p^a\),
\[
\phi(p^a)^2=p^{2a-2}(p-1)^2\ge p^a,
\]
and for \(2^a\),
\[
\phi(2^a)^2=2^{2a-2}\ge 2^{a-1}.
\]
Multiplying over prime powers gives
\[
\phi(m)^2\ge \frac m2
\]
for every \(m\ge 1\). Hence if \(\phi(m)=n\), then
\[
m\le 2n^2.
\]
Thus \(f_{\min}(n)\) and \(f_{\max}(n)\) are genuine finite quantities whenever \(n\in\phi(\N)\).

We shall use the following standard unconditional estimates: the prime number theorem in the form
\[
\vartheta(y):=\sum_{p\le y}\log p \sim y,
\]
Mertens' product theorem
\[
\prod_{p\le y}\left(1-\frac1p\right)^{-1}
\sim e^\gamma \log y,
\]
and Linnik's theorem, which says that there are absolute constants \(C,L>0\) such that, for every integer \(a\ge 1\), the least prime \(\ell\equiv 1\pmod a\) satisfies
\[
\ell\le C a^L.
\]
Increasing \(L\) if necessary, we may assume \(L\ge 1\).

We shall also need the following familiar consequence of the prime number theorem and Mertens' theorem.

\begin{lemma}
As \(T\to\infty\),
\[
\max_{1\le m\le T}\frac m{\phi(m)}
=(e^\gamma+o(1))\log\log T.
\]
\end{lemma}

\begin{proof}
Let \(p_1<p_2<\cdots\) be the sequence of primes and put
\[
N_k:=\prod_{i=1}^k p_i,
\qquad N_0:=1.
\]
Choose \(k=k(T)\) so that
\[
N_k\le T<N_{k+1}.
\]
If \(m\le T\) and \(r=\omega(m)\), then
\[
N_r\le \rad(m)\le m\le T,
\]
so \(r\le k\). Since the function \(p\mapsto p/(p-1)\) is decreasing on the primes,
\[
\frac m{\phi(m)}
=\prod_{p\mid m}\frac p{p-1}
\le \prod_{i=1}^r \frac{p_i}{p_i-1}
\le \prod_{i=1}^k\frac{p_i}{p_i-1}.
\]
Equality is attained at \(m=N_k\). Hence
\[
\max_{1\le m\le T}\frac m{\phi(m)}
=
\prod_{i=1}^k\frac{p_i}{p_i-1}.
\]
By Mertens' theorem,
\[
\prod_{i=1}^k\frac{p_i}{p_i-1}
=
(e^\gamma+o(1))\log p_k.
\]
On the other hand,
\[
\log N_k=\vartheta(p_k)\sim p_k,
\]
and since \(T<N_{k+1}\),
\[
\log T<\log N_{k+1}=\vartheta(p_{k+1})\sim p_{k+1}\sim p_k.
\]
Together with \(\log T\ge \log N_k\sim p_k\), this gives \(p_k\sim \log T\). Therefore
\[
\log p_k\sim \log\log T,
\]
and the lemma follows.
\end{proof}

\section{The asymptotic formula}

\begin{theorem}
As \(x\to\infty\),
\[
\mathcal R(x)
=
(e^\gamma+o(1))\log\log x .
\]
\end{theorem}

\begin{proof}
We prove the upper and lower bounds separately.

First let \(n\le x\) be a totient value, and put
\[
M=f_{\max}(n),
\qquad
m=f_{\min}(n).
\]
Since \(\phi(M)=\phi(m)=n\),
\[
\frac M m
=
\frac{M/\phi(M)}{m/\phi(m)}.
\]
Because \(m/\phi(m)\ge 1\), we have
\[
\frac M m
\le \frac M{\phi(M)}.
\]
Also, from the fibre-finiteness estimate above,
\[
M\le 2n^2\le 2x^2.
\]
Therefore
\[
\frac{f_{\max}(n)}{f_{\min}(n)}
\le
\max_{1\le t\le 2x^2}\frac t{\phi(t)}
=
(e^\gamma+o(1))\log\log(2x^2)
=
(e^\gamma+o(1))\log\log x.
\]
Taking the maximum over all totient values \(n\le x\) gives
\[
\mathcal R(x)
\le
(e^\gamma+o(1))\log\log x.
\]

It remains to construct totient fibres with this order of separation. Let \(y\to\infty\), and set
\[
P_y:=\prod_{p\le y}p,
\qquad
A_y:=\prod_{p\le y}(p-1).
\]
By Linnik's theorem, choose a prime \(\ell=\ell_y\) with
\[
\ell\equiv 1\pmod {A_y},
\qquad
\ell\le C A_y^L.
\]
Put
\[
U_y:=\frac{\ell-1}{A_y}.
\]
Let
\[
Q_y:=\prod_{\substack{q\mid U_y\\ q>y}}q,
\]
the squarefree product of the prime divisors of \(U_y\) exceeding \(y\). Define
\[
a_y:=\ell Q_y,
\qquad
b_y:=P_y U_y Q_y.
\]
We claim that \(\phi(a_y)=\phi(b_y)\).

Since \(U_y<\ell\), the prime \(\ell\) does not divide \(Q_y\), and hence
\[
\phi(a_y)
=
(\ell-1)\prod_{\substack{q\mid U_y\\ q>y}}(q-1)
=
A_yU_y\prod_{\substack{q\mid U_y\\ q>y}}(q-1).
\]
The prime divisors of \(b_y\) are exactly the primes \(p\le y\), together with the primes \(q>y\) dividing \(U_y\). Therefore
\[
\begin{aligned}
\phi(b_y)
&=
b_y
\prod_{p\le y}\left(1-\frac1p\right)
\prod_{\substack{q\mid U_y\\ q>y}}\left(1-\frac1q\right)  \\
&=
P_yU_yQ_y\cdot \frac{A_y}{P_y}
\cdot
\prod_{\substack{q\mid U_y\\ q>y}}\frac{q-1}{q} \\
&=
A_yU_y
\prod_{\substack{q\mid U_y\\ q>y}}(q-1).
\end{aligned}
\]
Thus
\[
\phi(a_y)=\phi(b_y)=:n_y.
\]

Their ratio is
\[
\frac{b_y}{a_y}
=
\frac{P_yU_yQ_y}{\ell Q_y}
=
\frac{P_y}{A_y}\cdot \frac{\ell-1}{\ell}.
\]
By Mertens' theorem,
\[
\frac{P_y}{A_y}
=
\prod_{p\le y}\frac p{p-1}
=
(e^\gamma+o(1))\log y.
\]
Also \(\ell\to\infty\), so
\[
\frac{\ell-1}{\ell}=1+o(1).
\]
Consequently
\[
\frac{b_y}{a_y}
=
(e^\gamma+o(1))\log y.
\]
For large \(y\), this ratio is greater than \(1\), and hence
\[
\frac{f_{\max}(n_y)}{f_{\min}(n_y)}
\ge
\frac{b_y}{a_y}
=
(e^\gamma+o(1))\log y.
\]

It remains to ensure that \(n_y\le x\) for a suitable choice of \(y\). Since \(Q_y\le U_y\),
\[
n_y
=
\phi(a_y)
=
A_yU_y
\prod_{\substack{q\mid U_y\\ q>y}}(q-1)
\le
A_yU_yQ_y
\le
A_yU_y^2.
\]
Using \(U_y=(\ell-1)/A_y\) and \(\ell\le C A_y^L\), we get
\[
n_y
\le
A_y\left(\frac{\ell-1}{A_y}\right)^2
\ll
A_y^{2L-1}.
\]
Moreover,
\[
\log A_y
=
\sum_{p\le y}\log(p-1)
=
\sum_{p\le y}\log p+\sum_{p\le y}\log\left(1-\frac1p\right)
=
\vartheta(y)+o(y)
=
(1+o(1))y.
\]
Hence
\[
\log n_y
\le
(2L-1+o(1))y.
\]

Now choose
\[
y=\frac{1}{4L}\log x.
\]
Then
\[
\log n_y
\le
\left(\frac{2L-1}{4L}+o(1)\right)\log x
<
\log x
\]
for all sufficiently large \(x\). Thus \(n_y\le x\), and so
\[
\mathcal R(x)
\ge
\frac{f_{\max}(n_y)}{f_{\min}(n_y)}
\ge
(e^\gamma+o(1))\log y.
\]
Since
\[
\log y
=
\log\log x+O(1),
\]
we conclude that
\[
\mathcal R(x)
\ge
(e^\gamma+o(1))\log\log x.
\]
Together with the upper bound, this proves
\[
\mathcal R(x)
=
(e^\gamma+o(1))\log\log x.
\]
\end{proof}

\section{A permanence observation}

The following elementary observation addresses the natural question of whether a single nontrivial totient collision forces infinitely many such collisions.

\begin{proposition}
Suppose \(a>b\) and \(\phi(a)=\phi(b)=n\). Then there are infinitely many distinct totient values \(N\) such that
\[
\frac{f_{\max}(N)}{f_{\min}(N)}\ge \frac ab.
\]
In particular, the existence of one nontrivial totient fibre implies the existence of infinitely many nontrivial totient fibres.
\end{proposition}

\begin{proof}
Choose any prime \(r\nmid ab\). Since \(r\) is coprime to both \(a\) and \(b\),
\[
\phi(ra)=(r-1)\phi(a)=(r-1)n
\]
and
\[
\phi(rb)=(r-1)\phi(b)=(r-1)n.
\]
Thus the totient value
\[
N_r:=(r-1)n
\]
has at least the two preimages \(ra\) and \(rb\), and
\[
\frac{ra}{rb}=\frac ab.
\]
Therefore
\[
\frac{f_{\max}(N_r)}{f_{\min}(N_r)}
\ge
\frac ab.
\]
As \(r\) ranges over the infinitely many primes not dividing \(ab\), the integers \(N_r=(r-1)n\) are distinct. Hence infinitely many such totient values exist.
\end{proof}

\end{document}
